% from ega4-11-fr.tex
\oldpage[IV-3]{162}
en général, que la famille $(\varphi_\lambda)$ soit séparante. Par exemple,
supposons que $A$ soit un anneau de valuation discrète, d'idéal maximal
$\mathfrak m$, et considérons la famille des homomorphismes
$\varphi_k:A\to A/\mathfrak m^k$, dont l'intersection des noyaux
$\mathfrak m^k$ est réduite à 0~; cette famille n'est toutefois pas séparante,
car les fibres de tous les $\mathfrak m^k$ au point générique $x$ de
$X=\operatorname{Spec}(A)$ (qui est ouvert dans $X$) sont égales à
$\mathcal O_x=k(x)$, corps des fractions de $A$, et leur intersection n'est
donc pas réduite à 0.
\end{env}

\begin{env}[11.9.5]
\label{IV.11.9.5-fr}
Nous allons principalement nous occuper dans ce qui suit du problème du
\emph{changement de base} pour les familles séparantes. Les notations étant
celles de (11.9.1), considérons un morphisme $g:X'\to X$ et posons
$\mathcal F'=\mathcal F\otimes_{\mathcal O_X}\mathcal O_{X'}
=g^*(\mathcal F)$, et, pour tout $\lambda$,
$Z'_\lambda=Z_\lambda\times_X X'$, $f'_\lambda=f_\lambda\times 1$,
$\mathcal G'_\lambda=\mathcal G_\lambda\otimes_{\mathcal O_{Z_\lambda}}
\mathcal O_{Z'_\lambda}=(g'_\lambda)^*(\mathcal G_\lambda)$, où
$g'_\lambda:Z'_\lambda\to Z_\lambda$ est la projection canonique. Pour tout
$\lambda$, on désigne alors par
$u'_\lambda:\mathcal F'\to(f'_\lambda)_*(\mathcal G'_\lambda)$
l'homomorphisme obtenu comme suit~: soit
\[
\psi:(f_\lambda)_*(\mathcal G_\lambda)\longrightarrow
(f_\lambda)_*((g'_\lambda)_*((g'_\lambda)^*(\mathcal G_\lambda)))
=g_*((f'_\lambda)_*(\mathcal G'_\lambda))
\]
l'homomorphisme $(f_\lambda)_*(\rho_{\mathcal G_\lambda})$, où
$\rho_{\mathcal G_\lambda}$ est l'homomorphisme canonique
$(\mathbf 0_{\mathrm I},4.4.3.2)$ correspondant à $g'_\lambda$. Alors
$u'_\lambda$ est défini comme le composé
\[
g^*(\mathcal F)\xrightarrow{g^*(u_\lambda)}
g^*((f_\lambda)_*(\mathcal G_\lambda))
\xrightarrow{g^*(\psi)}g^*g_*((f'_\lambda)_*(\mathcal G'_\lambda))
\xrightarrow{\sigma}(f'_\lambda)_*(\mathcal G'_\lambda)
\]
où $\sigma=\sigma_{(f'_\lambda)_*(\mathcal G'_\lambda)}$ est
l'homomorphisme canonique $(\mathbf 0_{\mathrm I},4.4.3.3)$ correspondant à
$g$. On dira pour abréger que $u'_\lambda$ est \emph{déduit de $u_\lambda$
par le changement de base} $g$. Lorsque $f_\lambda:Z_\lambda\to X$ est un
morphisme \emph{affine}, on a
$(f'_\lambda)_*(\mathcal G'_\lambda)=g^*((f_\lambda)_*(\mathcal G_\lambda))$,
et dans ce cas on a donc simplement $u'_\lambda=g^*(u_\lambda)$.

On peut aussi interpréter $u'_\lambda$ de la façon suivante~: il suffit de
connaître la valeur de $u'_\lambda(t')$, lorsque $t'$ est une section de
$\mathcal F'$ au-dessus d'un ouvert $U'$ de $X'$, du type particulier
suivant~: $t'$ est la restriction à $U'$ de l'image canonique par
$\Gamma(U,\mathcal F)\to\Gamma(g^{-1}(U),\mathcal F')$ d'une section $t$ de
$\mathcal F$ au-dessus d'un ouvert $U$ de $X$ contenant $g(U')$ (ces sections
engendrant en effet le $\mathcal O_{X'}$-Module $\mathcal F'$
$(\mathbf 0_{\mathrm I},3.7.1)$). Considérons la section $u_\lambda(t)$ de
$\mathcal G_\lambda$ au-dessus de $f_\lambda^{-1}(U)$, et son image canonique
$t''$ par
$\Gamma(f_\lambda^{-1}(U),\mathcal G_\lambda)\to
\Gamma((g'_\lambda)^{-1}(f_\lambda^{-1}(U)),\mathcal G'_\lambda)$~; alors
$u'_\lambda(t')$ est la restriction de $t''$ à $(f'_\lambda)^{-1}(U')$.
Considérons en particulier le cas où $Z_\lambda$ est un sous-préschéma induit
sur un ouvert de $X$, $\mathcal G_\lambda=j_\lambda^*(\mathcal F)$, où
$j_\lambda:Z_\lambda\to X$ est l'injection canonique, et où $u_\lambda$ est
l'homomorphisme canonique
$\rho_{\mathcal F}:\mathcal F\to(j_\lambda)_*(j_\lambda^*(\mathcal F))
=(j_\lambda)_*(\mathcal G_\lambda)$ $(\mathbf 0_{\mathrm I},4.4.3.2)$
correspondant à $j_\lambda$. Alors $Z'_\lambda$ est induit sur un ouvert de
$X'$, et l'interprétation précédente montre que $u'_\lambda$ n'est autre que
l'homomorphisme canonique
$\rho_{\mathcal F'}:\mathcal F'\to
(j'_\lambda)_*((j'_\lambda)^*(\mathcal F'))
=(j'_\lambda)_*(\mathcal G'_\lambda)$ correspondant à l'injection canonique
$j'_\lambda:Z'_\lambda\to X'$.
\end{env}

\begin{env}[11.9.6]
\label{IV.11.9.6-fr}
Sous les conditions de (11.9.5), supposons que $X$ et $X'$ soient affines, et
que l'on veuille prouver que pour toute section $t'$ de $\mathcal F'$ au-dessus
de $X'$, dont les images par tous les $u'_\lambda$ sont nulles, alors $t'$ est
elle-même nulle. Alors, on peut aussi se borner au cas où $Z_\lambda=X$ pour
tout $\lambda\in L$. En effet, avec les notations de (11.9.3) et de (11.9.5),
si l'on pose $Y'_{\lambda\mu}=(g'_\lambda)^{-1}(Y_{\lambda\mu})$,
l'homomorphisme $v'_{\lambda\mu}$ déduit de $v_{\lambda\mu}$ par le changement
de base $g$ n'est autre que l'homomorphisme canonique
